\section*{5 Unordered Multiselection Over Large Alphabets (author source pp.11:9--11:13)}
In this section we construct a circuit that implements a relaxation of multiselection. Informally, this relaxation allows the output elements to appear in any order (and possibly with repetitions), rather than the same order in which they were queried. Additionally, rather than selecting individual bits (which is our eventual goal)---we consider a generalization to a larger alphabet size.
\begin{originalresult}{Definition 12}
Unordered multiselection (with $q$ queries, over alphabet $\Sigma$) is the search problem defined by the relation
\[\uSel_\Sigma^{n\to q}\stackrel{\rm def}=\{((x,i),y)\in(\Sigma^n\times[n]^q)\times(\Sigma\cup\{\bot\})^q:\{y_k:y_k\ne\bot\}_{k\in[q]}=\{x_{i_k}\}_{k\in[q]}\}.\]
\end{originalresult}
\subsection*{5.1 Superconcentrators and Routing}
We recall the notion of a superconcentrator, which is the main technical tool involved in our construction of Boolean circuits for unordered multiselection.
\begin{originalresult}{Definition 13 (Networks)}
A network $N$ is a tuple $N=(V,E,A,B)$, where $G=(V,E)$ is a directed acyclic graph and $A,B\subseteq V$ are respectively sets of sources and sinks in $G$.
If $|A|=|B|=n$, then $N$ is said to be an $n$-network. The size of $N$, denoted by $|N|$, is defined to be $|E|$; the degree of $N$ is the degree of $G$; and the depth of $N$ is the longest path length in $G$.
A family of $n$-networks $\{N_n\}_{n\in\ZZ^+}$ is said to be log-space uniform if there an algorithm that on input $1^n$ runs in $O(\log n)$ space and outputs $G$ (say with an adjacency list representation) and $A,B$ (say represented by indicator strings).
\end{originalresult}
\begin{originalresult}{Definition 14 (Superconcentrators)}
An $n$-network $N=(V,E,A,B)$ is said to be a superconcentrator if for all $q\in[n]$ and all subsets $X\subseteq A$ and $Y\subseteq B$ with $|X|=|Y|=q$, there exist $q$ vertex-disjoint paths from $X$ to $Y$.
\end{originalresult}
\begin{originalresult}{Definition 15}
Let $S=(V,E,A,B)$ be an $n$-superconcentrator. The routing problem for $S$ is a search problem $\operatorname{Route}_S$ whose input is a pair of (indicator strings for) sets $X\subseteq A,Y\subseteq B$ with $|X|=|Y|$. A valid corresponding output is any (indicator string for a) set $R\subseteq E$ of edges such that $R$ is a union of $q$ vertex-disjoint paths from $X$ to $Y$, where $q=|X|=|Y|$.
\end{originalresult}
Note that whenever $S$ is an $n$-superconcentrator, the search problem $\operatorname{Route}_S$ is total. Naturally, it is desirable for a superconcentrator $S$ to admit efficient algorithms for $\operatorname{Route}_S$. The state of the art in this respect is Pippenger's construction of a ``self-routing'' superconcentrator.
Informally, a delay-$d$ self-routing protocol for a superconcentrator $(V,E,A,B)$ associates each vertex $v\in V$ with a finite automaton that can communicate synchronously with the automata on neighboring vertices (where $u$ is said to neighbor $v$ if $(u,v)$ or $(v,u)$ is in $E$). For any sets $X\subseteq A$ and $Y\subseteq B$ with $|X|=|Y|=q$, if the automata in $X$ and $Y$ are initialized with state 1, and all other automata are initialized with state 0, then in $d$ steps the automata jointly compute $q$ vertex-disjoint paths from $X$ to $Y$ by on the $d$th step transmitting 1 on all edges in those paths and 0 on other edges.
\begin{originalresult}{Definition 16}
A self-routing protocol with delay $d$ for an $n$-superconcentrator $S=(V,E,A,B)$ is a tuple $(\Sigma,\delta)$, where:
\begin{itemize}
\item $\Sigma$ is a finite ``alphabet'' set that contains $\{0,1\}$; and
\item $\delta:\Sigma^{V\cup E}\to\Sigma^{V\cup E}$ is a function such that:
\begin{itemize}
\item for every vertex $v\in V$ with incident edges $E_v$, $\delta(q)_v$ depends as a function of $q$ only on $(q_e)_{e\in E_v}$.
\item for every edge $e=(u,v)\in E$, $\delta(q)_e$ depends as a function of $q$ only on $q_u$ and $q_v$.
\end{itemize}
\end{itemize}
For any sets $X\subseteq A$ and $Y\subseteq B$ with $|X|=|Y|$, if we define $q^{(0)}\in\Sigma^{V\cup E}$ by
\[q_v^{(0)}\stackrel{\rm def}=\begin{cases}1&\text{if }v\in X\cup Y,\\0&\text{if }v\in V\setminus(X\cup Y),\end{cases}\]
and $q_e^{(0)}\stackrel{\rm def}=0$ for all $e\in E$, then the set
\[R\stackrel{\rm def}=\{e\in E:\delta^d(q^{(0)})_e=1\}\]
satisfies $((1_X,1_Y),1_R)\in\operatorname{Route}_S$.
Such a self-routing protocol for an $n$-superconcentrator is said to be log-space uniform if there is a log-space Turing machine that on input $1^n$ outputs:
\begin{itemize}
\item a description of $\Sigma$;
\item for each $v\in V$ with incident edges $E_v$, the truth table for $\delta(q)_v$ as a function of $(q_e)_{e\in E_v}$.
\item for each $e=(u,v)\in E$, the truth table for $\delta(q)_e$ as a function of $(q_u,q_v)$.
\end{itemize}
\end{originalresult}
\begin{originalresult}{Definition 17}
A self-routing superconcentrator is a superconcentrator $S$ with an associated self-routing protocol $(\Sigma,\delta)$ for $S$. A self-routing superconcentrator is said to be log-space uniform if both the superconcentrator and the associated self-routing protocol are log-space uniform.
\end{originalresult}
\begin{originalresult}{Theorem 18 ([25])}
There exists a log-space uniform self-routing $n$-superconcentrator $S$ with size $O(n)$, degree $O(1)$, and depth $O(\log n)$, where the self-routing protocol has alphabet size $O(1)$ and delay $O(\log n)$.
\end{originalresult}
\begin{originalresult}{Corollary 19}
There exists a log-space uniform $n$-superconcentrator $S$ with size $O(n)$, degree $O(1)$, and depth $O(\log n)$, such that $\operatorname{Route}_S$ is solvable by a log-space uniform Boolean circuit of size $O(n\log n)$ and depth $O(\log n)$.
\end{originalresult}
\begin{proof}
Let $S=(V,E,A,B)$ be the log-space uniform self-routing superconcentrator given by Theorem 18, let the self-routing protocol be given by $(\Sigma,\delta)$, and let $d=O(\log n)$ denote the delay of this self-routing protocol. To prove the corollary, we need to construct a log-space uniform Boolean circuit for solving $\operatorname{Route}_S$, with size $O(n\log n)$ and depth $O(\log n)$.
We are given $1_X$ and $1_Y$. The circuit starts by computing $q^{(0)}\in\Sigma^{V\cup E}$ as defined in Definition 16. By construction this can be done with a circuitry of size $O(n)$ and depth $O(1)$. The circuit next computes $\delta^d(q^{(0)})$. Recall that each output symbol of $\delta(q)$ depends on at most $\deg(S)=O(1)$ symbols of $q$. Since the alphabet $\Sigma$ is also constant-sized, this means that $\delta$ is implementable by a (log-space uniform) circuit of size $O(|V\cup E|)=O(n)$ and depth $O(1)$. Iterating this circuit $d$ times, we compute $\delta^d(q^{(0)})$ with log-space uniform circuitry of size $O(d\cdot n)$ and depth $O(d)$. Finally, we extract from $\delta^d(q^{(0)})$ the indicator string $1_R$ for the set
\[R\stackrel{\rm def}=\{e\in E:\delta^d(q^{(0)})_e=1\}.\]
This takes (log-space uniform) circuitry of size $O(|E|)=O(n)$ and depth $O(1)$.
In total, this circuit for computing $1_R$ is log-space uniform and has size $O(n\cdot d)=O(n\log n)$ and depth $O(d)=O(\log n)$. By the definition of a self-routing protocol, $R$ satisfies $((1_X,1_Y),1_R)\in\operatorname{Route}_S$, i.e. the circuit solves $\operatorname{Route}_S$.
\end{proof}
\subsection*{5.2 From Superconcentrators To Unordered Multiselection}
\begin{originalresult}{Proposition 20}
Let $n\in\ZZ^+$, let $s=s(n)$ satisfy $s=\Omega(\log^2n)$, and let $\Sigma=\bits^s$. For any $q=q(n)\in[n]$ there exists a log-space uniform Boolean circuit for $\uSel_\Sigma^{n\to q}$ with size $O(n\cdot s)$ and depth $O(\log n)$.
\end{originalresult}
\begin{proof}
Let $S=(V,E,A,B)$ be a log-space uniform $n$-superconcentrator with size $O(n)$, depth $O(\log n)$, degree $O(1)$ such that $\operatorname{Route}_S$ is solvable by a log-space uniform Boolean circuit of size $O(n\log n)$ and depth $O(\log n)$. (Such an $S$ is guaranteed to exist by Corollary 19). Order the elements of $A$ and $B$ arbitrarily so that $A=\{a_1,\ldots,a_n\}$ and $B=\{b_1,\ldots,b_n\}$.
\paragraph{The Circuit.} Our circuit, taking as input $x\in\Sigma^n$ and $(i_1,\ldots,i_q)\in[n]^q$, is constructed in two parts.
The first part of our circuit computes indicator strings $1_X\in\bits^A$, $1_Y\in\bits^B$, and $1_R\in\bits^E$ for sets $X\subseteq A,Y\subseteq B$, and $R\subseteq E$ defined as follows. The set $X$ is defined as $\{a_{i_1},\ldots,a_{i_q}\}$, $Y$ is defined as $\{b_1,\ldots,b_{q'}\}$, where $q'$ is the cardinality of $X$ (which may be smaller than $q$ due to repetitions), and $R$ is the edges of $q'$ vertex-disjoint paths from $X$ to $Y$.
The second part of our circuit consists of a gadget $g_v$ for each vertex $v\in V$. For a vertex $v\in A$, say $v=a_j$, the gadget $g_v$ simply outputs $x_j$ if $v\in X$ and $\bot$ otherwise. For other $v$, i.e. with positive in-degree $d$, the gadget $g_v$ is a circuit for $\Sel_{\Sigma\cup\{\bot\}}^{d\to1}$ (e.g. the circuit from Proposition 2). To construct the inputs for $g_v$, first order the in-neighbors of $v$ arbitrarily as $u_1,\ldots,u_d$. The data input for $g_v$ is then constructed by taking the $i$th symbol, for any $i\in[d]$, to be the output of $g_{u_i}$. The selector input for $g_v$ is constructed to be (the unique) $i^\star\in[d]$ such that $(u_{i^\star},v)\in R$ if such an $i^\star$ exists, and arbitrary otherwise.
Finally, the outputs of our circuit are the outputs of $g_{b_1},\ldots,g_{b_q}$.
\paragraph{Size and Depth.} In the first part of our circuit, the computation of $1_X$ uses log-space uniform circuitry of size $O(n\log^2n)=O(n\cdot s)$ and depth $O(\log n)$ by Lemma 21 below. Next, $1_Y$ is obtained by sorting $1_X$ with log-space uniform circuitry of size $O(n\log n)$ and depth $O(\log n)$ (this is possible by Lemma 8). Finally, the computation of $1_R$ uses log-space uniform circuitry of size $O(n\log n)$ and depth $O(\log n)$ because of our choice of superconcentrator $S$. In total the circuitry of this part has size $O(n\cdot s)$ and depth $O(\log n)$.
In the second part of our circuit, there are $O(n)$ gadgets (because $S$ has size $O(n)$), and each gadget has size $O(s)$ (because $S$ has constant degree). Additionally for each non-source gadget there is constant-sized circuitry to compute the selector input as a function of $1_R$ (again because $S$ has constant degree). In total the circuitry size is $O(n\cdot s)$. As for depth, each gadget has constant depth. Since the output of gadget $u$ is an input to gadget $v$ iff $(u,v)$ is an element of $E$, and $S$ has depth $O(\log n)$, the circuitry here also has depth $O(\log n)$.
\paragraph{Correctness.} The main claim that we use to establish correctness is that for every $a_i\in X$ and every vertex $v\in V$, if $v$ lies on a path in $R$ from $a_i$ to $B$, then the output of $g_v$ is $x_i$, and otherwise the output of $g_v$ is $\bot$. Here $a_i$ is well-defined as a partial function of $v$ because $R$ consists only of vertex-disjoint paths. This claim suffices for correctness because there are paths in $R$ from each $a_{i_k}$ to $B$, but not from any $a\in A\setminus\{i_1,\ldots,i_q\}$, so the set of non-$\bot$ outputs is exactly $\{x_{i_1},\ldots,x_{i_q}\}$.
We prove the claim by induction on the depth of $v$ (i.e. the maximum length of a path ending in $v$). The base case is when $v$ has depth 0, i.e. $v$ is a source vertex $a_i$ in $(V,E)$. In this case we defined $g_v$ to output $x_i$.
For the inductive step, consider a vertex $v$ with positive depth, and suppose that the claim holds for all $u$ of lesser depth, and in particular for the in-neighbors $u_1,\ldots,u_d$ of $v$. Now if $v$ lies on a path $p$ from some $a_i$ to $Y$, this path must go through $u_j$ for some $j\in[d]$, and then the inductive hypothesis implies that the output of $g_{u_j}$ is $x_i$. The construction of the selector input for $g_v$ ensures that $g_v$ also outputs $x_i$, which completes the proof of the claim.
\end{proof}
\subsection*{5.3 Computing Set Indicator Strings}
\begin{originalresult}{Lemma 21}
For $q=q(n)\in\ZZ^+$, there is a log-space uniform Boolean circuit with size $O((n+q)\cdot\log^2(n+q))$ and depth $O(\log(n+q))$ that maps $(i_1,\ldots,i_q)\in[n]^q$ to the indicator string $1_{\{i_1,\ldots,i_q\}}\in\bits^n$.
\end{originalresult}
\begin{proof}
On input $(i_1,\ldots,i_q)$, the circuit performs the following steps:
\begin{enumerate}
\item Construct the list $(1,\ldots,n,i_1,\ldots,i_q)\in[n]^{n+q}$, and sort it to obtain a non-decreasing list $(j_i)_{i\in[n+q]}$ with the property that $j_i$ appears more than once in the list iff $j_i\in\{i_1,\ldots,i_q\}$. This can be done with circuitry of size $O((n+q)\cdot\log^2(n+q))$ and depth $O(\log(n+q))$ by Lemma 6.
\item Label $j_i$ with 0 if the value $j_i$ appears once in the list. If $j_i$ appears more than once, we label the first occurrence with 1 and all other occurrences with $\bot$. This labels $j_i$ with 1 only if $j_i$ is in $\{i_1,\ldots,i_q\}$, and labels $j_i$ with 0 only if it isn't. More explicitly, we compute labels
\[b_i=\begin{cases}\bot&\text{if }i>1\text{ and }j_{i-1}=j_i,\\1&\text{otherwise, if }i<n+q\text{ and }j_i=j_{i+1},\\0&\text{otherwise}\end{cases}\]
for $i\in[n+q]$. This can be done with circuitry of size $O((n+q)\log n)$ and depth $O(\log\log n)$: first compute in parallel for each $i\in[n+q-1]$ whether or not $j_i=j_{i+1}$ (this circuit has size $O(\log n)$ and depth $O(\log\log n)$), and then compute each $b_i$ with a constant-sized circuit.
\item Sort the list $(j_i,b_i)_{i\in[n+q]}$ in order of increasing $j_i$, except that if $b_i=\bot$ we treat $j_i$ as $+\infty$. That is, we prepend each $j_i$ with 1 if $b_i=\bot$ and with 0 otherwise. We then take the first $n$ elements of the result to obtain a list $(i,b_i')_{i\in[n]}$ such that $b_i'\in\{0,1\}$ satisfies
\[b_i'=\begin{cases}1&\text{if }i\in\{i_1,\ldots,i_q\},\\0&\text{otherwise}.\end{cases}\]
In other words, $(b_1',\ldots,b_n')$ is $1_{\{i_1,\ldots,i_q\}}$, which is our desired output. This step can be done with circuitry of size $O((n+q)\cdot\log^2(n+q))$ and depth $O(\log(n+q))$ by Lemma 6.
\end{enumerate}
\end{proof}
